\textbf{COLUNI 2021 -- Questão 01}\par\medskip Leia o texto:

\begin{center}
\textit{Sonho Imigrante}

A terra do sonho é distante

e seu nome é Brasil

plantarei a minha vida

debaixo de céu anil.

Minha Itália, Alemanha

Minha Espanha, Portugal

talvez nunca mais eu veja

minha terra natal.

Aqui sou povo sofrido

lá eu serei fazendeiro

terei gado, terei sol

O mar de lá é tão lindo

natureza generosa

que faz nascer sem espinho

o milagre da rosa.

O frio não é muito frio

nem o calor é muito quente

e falam que quem lá vive

é maravilha de gente.
\end{center}

BRANT, F., Nascimento, M. \textbf{Sonho Imigrante}. Disponível em: <http://construtor.aprendebrasil.com.br/up121780001/7910605/Sonho%2520Imigrante.pdf>. Acesso em: 10 dez. 2020. (Adaptado)

\par

No livro \textit{A menina que fez a América}, de Ilka B. Laurito, e no poema \textit{Sonho Imigrante}, de Fernando Brant e de Milton Nascimento, é INCORRETO afirmar que retratam, em síntese:
\begin{enumerate}
\item o ponto de vista do imigrante.
\item o sonho de vida farta e generosa.
\item a situação de poucos europeus.
\item o desejo de encontrar o novo.
\end{enumerate}
